\section{Keep both signs in the fourth-moment charge}

Let $X$ be centered, with $\E X^2=1$ and $|X|\le R$ almost surely.
Write
\[
 \rho=\E|X|^3,\quad\beta=\E X|X|,\quad\gamma=\E X^3,
 \qquad K_1=\E[X|X|^3]-3\beta-\tfrac32\rho\gamma.
\]
There is no central-gap or small-noise hypothesis. Atoms, including an
atom at zero, are allowed. The aim is to bound the signed fourth-moment
charge $K_1$ without combining incompatible extremal signs.

Put $m=\E X_+=\E X_->0$ and
$a_\pm=\E X_\pm^2$, $b_\pm=\E X_\pm^3$, $c_\pm=\E X_\pm^4$.
Both $a_\pm$ are positive, and signwise Cauchy--Schwarz gives
\begin{equation}\label{signed:psd}
 a_\pm^2\le m b_\pm,\qquad b_\pm^2\le a_\pm c_\pm,
 \qquad c_\pm\le Rb_\pm,\qquad 2m\le1.
\end{equation}
The last bound follows from $\E|X|\le\sqrt{\E X^2}=1$.
These are inequalities for one law, not a sufficiency claim for a moment
relaxation.

\begin{lemma}[A signed one-variable enclosure]\label{signed:lemma}
Set
\[
 C(r)=Rr-\tfrac32r^2-3,\qquad a(r)=\tfrac32r-R/2.
\]
For either $\epsilon\in\{-1,1\}$, let
$h=1-\epsilon\beta$ and $\delta=\rho-\epsilon\gamma$. Then
\begin{equation}\label{signed:core}
 h>0,\qquad\delta\ge h^2,\qquad
 \epsilon K_1\le C(\rho)+3h+a(\rho)\delta-\frac{\delta^2}{2h}.
\end{equation}
\end{lemma}
\begin{proof}
Take $\epsilon=1$. Then $a_-=h/2$ and $b_-=\delta/2$.
The first and last inequalities of \eqref{signed:psd} yield
$h^2/4\le m\delta/2\le\delta/4$.
The other two inequalities yield
\[
 c_-\ge\frac{\delta^2}{2h},\qquad
 c_+\le\frac R2(\rho+\gamma).
\]
Insert these bounds into
$K_1=c_+-c_--3\beta-3\rho\gamma/2$, and substitute
$\beta=1-h$, $\gamma=\rho-\delta$. This gives
\eqref{signed:core}. Apply the same argument to $-X$ for
$\epsilon=-1$. Its $K_1$ is $-K_1(X)$, while $\rho$ and $R$ are unchanged.
\end{proof}

For clarity, the absolute sign-moment cap also has a short direct proof.
If $p_\pm=\mathbb P(\pm X>0)$, Cauchy gives
$m^2/a_\pm\le p_\pm$. Since $p_++p_-\le1$, it follows that
$m^2\le a_+a_-=(1-\beta^2)/4$.
Adding $a_\pm^2\le mb_\pm$ therefore gives
\begin{equation}\label{signed:beta}
 \rho\ge\frac{1+\beta^2}{\sqrt{1-\beta^2}}.
\end{equation}
In particular, a rational $b\in[0,1)$ is a valid bound for $|\beta|$
throughout $\rho\le u$ whenever
$(1+b^2)^2\ge u^2(1-b^2)$. The right side of
\eqref{signed:beta} is strictly increasing as a function of $|\beta|$.

\section{A closed-cell bound with three scalar gates}

Assume $\rho\in[l,u]$ and $|\beta|\le b<1$, and put $h_*=1+b$.
Ignoring the constraint $\delta\ge h^2$ and completing the square in
\eqref{signed:core} gives the universally valid preliminary bound
\begin{equation}\label{signed:easy}
 |K_1|\le B_R(\rho,b):=
 R\rho-\tfrac32\rho^2+3b+
 \frac{1+b}{2}(\tfrac32\rho-R/2)^2.
\end{equation}
Indeed its remaining coefficient $3+a(\rho)^2/2$ is positive, so
$h\le h_*$ can be inserted after square completion.
The maximum of the quadratic $B_R(r,b)$ on $[l,u]$ is at an endpoint,
or at its vertex if the quadratic is concave and the vertex lies inside.
Thus \eqref{signed:easy} already has a finite exact acceptance rule.

The constraint $\delta\ge h^2$ gives a stronger bound.

\begin{theorem}\label{signed:theorem}
Suppose, in addition, that
\begin{align}
 h_*&\ge a(u),\label{signed:gate1}\\
 3+2a(l)h_*-\tfrac32h_*^2&\ge0,\label{signed:gate2}\\
 R-3u+\tfrac32h_*^2&\ge0.\label{signed:gate3}
\end{align}
Then
\begin{equation}\label{signed:strong}
 |K_1|\le S(R,u,b):=
 Ru-\tfrac32u^2-3+3h_*+
 (\tfrac32u-R/2)h_*^2-\tfrac12h_*^3.
\end{equation}
All cell endpoints are included.
\end{theorem}
\begin{proof}
For fixed $a=a(\rho)$, maximizing the right side of
\eqref{signed:core} over $\delta\ge h^2$ leaves $C(\rho)+\psi_a(h)$,
where
\[
 \psi_a(h)=
 \begin{cases}
 h(3+a^2/2),&0<h\le a,\\
 3h+ah^2-h^3/2,&h\ge a.
 \end{cases}
\]
The first branch exists only if $a>0$ and is increasing.
The second branch has derivative $3+2ah-3h^2/2$, a concave quadratic
in $h$. On $[0,h_*]$ its endpoint values are $3$ and at least the
left side of \eqref{signed:gate2}, because $a\ge a(l)$.
It is therefore nonnegative throughout. The two branches agree at
$h=a$, so $\psi_a$ is increasing up to $h_*$.
By \eqref{signed:gate1}, the value at $h_*$ is the second branch.
The resulting expression
$C(\rho)+3h_*+a(\rho)h_*^2-h_*^3/2$ has derivative
$R-3\rho+3h_*^2/2$, which is nonnegative on $[l,u]$ by
\eqref{signed:gate3}. Insert $\rho=u$ to obtain \eqref{signed:strong}.
\end{proof}

For the last recorded initial cell, the inputs are
\[
 R=96/25,\quad l=173303/96000,\quad u=7361/4000,
 \quad b=32078117/50000000.
\]
The three gate margins are positive; the derivative margin in
\eqref{signed:gate2} exceeds $1.5445$.
Exact rational substitution gives
\[
 S(R,u,b)<\frac{3964267}{10^6}<\frac{397}{100}.
\]
This improves the retained bound $|K_1|\le5.8704858543$ on this cell.
The same gates pass on cells $10$--$23$ of the recorded 24-cell cover
(and also cell $0$, where the retained bound is better).
Where a gate fails, use \eqref{signed:easy} or retain the earlier bound;
the stronger formula is not extrapolated through a failed gate.

\section{A short fifth-moment square and its unsigned use}

Let $M_4=\E X^4$ and $M_5=\E|X|^5$. The identity
\[
 0\le\E\{|X|^3(|X|-\rho)^2\}
 =M_5-2\rho M_4+\rho^3
\]
proves
\begin{equation}\label{signed:fifth}
 \frac\rho2M_4-\frac14M_5\le\frac{\rho^3}{4}.
\end{equation}
This is sharp for the standardized symmetric three-point law on
$\{-\rho,0,\rho\}$ with masses $1/(2\rho^2)$ at each nonzero point,
for every $\rho\ge1$. A zero atom is therefore important in testing it.
This inequality is already used in the retained signed code. The present
step uses it in the earlier unsigned coefficient, replacing the weaker
$8\rho^3/27$ term; no novelty claim is attached to the inequality.

Let $M\ge M_4$ and let $\Gamma(R,l)\ge0$ be the retained negative-gap
coefficient. In the retained coarse constant envelope, write
\[
 A_1=\tfrac32b(2u+1)+3\Gamma(R,l)u
       +\tfrac38u^2(u+1)+\tfrac8{27}u^3-l/2,
 \quad A_2=\tfrac32b(2u+1)+2uM.
\]
The retained second branch is $A_2-l/2$. With exactly the same gap and
moment premises, the new constant is
\begin{equation}\label{signed:A}
 A_{\rm new}=\min\{A_1-5u^3/108,\ A_2-l/2\}.
\end{equation}
In particular its last-cell value is less than $10.406486$, compared
with $10.6950063919$ previously. This step retains every other charge;
it is not a new proof of the gap estimate represented by $\Gamma$.

\section{Checks and limits}

The exact checker reconstructs the 24 closed $\rho$ cells, their rational
moment caps and the retained scalar comparison formulas; it imports no
archived program. It checks \eqref{signed:beta}, all three gates where
\eqref{signed:strong} is used, and the endpoint/vertex comparisons for
\eqref{signed:easy}. Failed monotonicity gates are kept as failures.
Separate output columns distinguish the new universal bound from its
minimum with the earlier conditional cap. No Fourier integral is
evaluated by this check.

Actual-law controls include symmetric and skew Bernoulli laws, the sharp
three-point example for \eqref{signed:fifth}, and the centered,
variance-one law with atoms and masses
\[
 (-96/25,-1/2,37/25),\qquad
 (325/22211,11708/16533,3650/13167).
\]
For the last law,
\[
 \rho=29884193/16458750\in[l,u],\quad
 \beta=70573/329175>0,\qquad\gamma=-23/1250<0.
\]
Thus $\beta\gamma\ge0$ is false even inside the last acquired moment
cell. Also
\[
 K_1=-102409783933/41146875000<0,
\]
so $K_1\gamma\le0$ is false there. Neither tempting sign shortcut is used.

The conclusions above are moment inequalities under an actual support
bound. A raw-remainder or conserved-resource application must still
supply the fixed-target contact/transfer identity, the acquired support,
and the full low- and high-frequency estimates. No initial-acquisition
theorem, local exclusion, saved certificate job, or global bound follows
from these coefficient improvements alone.
